\chapter{Computation and experimental design}\label{sec:methods}
\section{Independent implementations and validation}\label{sec:validation}
A reference simulation scans the grid and performs one elementary toppling at a time. The faster algorithm keeps a queue of unstable sites and performs $\lfloor h(x)/4\rfloor$ consecutive legal topplings at each visit, adding that count to the odometer. The queue holds at most one entry per site. This avoids repeated scans of stable sites and duplicate queue entries. We compared both Python and C versions with the reference.

Every event was checked against the mass balance. We also checked actual grid mass using height frequencies every 128 added grains and direct sums during each run. A separate array-based simulation performed this check after every validation event. Heights and counts use signed 64-bit integers, and repeated batch locations receive every grain assigned to them.

The comparisons cover all $256\times4=1,024$ stable $2\times2$ additions and 1,000 bounded random configurations for each implementation. Forward, reverse and random legal toppling orders give identical results. Tests also cover rectangular and degenerate grids, repeated topplings, all four observables and the corrected handout example. Extension checks verify periodic connections, dissipation and distances on the cylinder, biased-input probabilities and frequencies, and batch odometer sums and area unions. Exact enumeration gives 192 recurrent $2\times2$ states and verifies the uniform invariant measure under uniform and unequal positive input probabilities, including four successive additions.

Known discrete distributions, sparse samples and autoregressive series were used to check the statistical calculations. For each fit, we compared several initial parameter choices and checked the best converged solution against an independent projected-gradient calculation. All 144 aggregate fits passed these checks.

\section{Profiling and performance}\label{sec:performance}
Profiling identified repeated neighbour construction as a major cost, so we stored the neighbour lists. Profiling at $L=64$ then showed that processing the active-site queue took most of the time. We implemented this algorithm in C, with Python generating the input sites, and calculated observables as sites toppled. Table~\ref{tab:performance} compares stabilization and observable calculation for a fixed input sequence.

\begin{table}[htbp]
\centering\small
\begin{tabular}{@{}lrrr@{}}
\toprule
Implementation & Median (ms) & Range (ms) & Speed ratio \\
\midrule
Elementary scan, Python & 552.216 & [531.345, 562.245] & 1.0 \\
Bulk queue, Python & 107.030 & [106.292, 107.353] & 5.2 \\
C queue, one event per call & 95.477 & [94.000, 107.416] & 5.8 \\
C queue, full input sequence & 0.167 & [0.166, 0.205] & 3298.7 \\
\bottomrule
\end{tabular}

\caption{Stabilization times for the same $L=16$ calculation: \R{benchmarkEvents} predetermined inputs from an all-three grid, giving \R{benchmarkTopples} elementary topplings. Each implementation first completes one untimed run. Each of five timing samples then repeats complete calculations for at least 0.2\,s, resetting the grid each time. Reported times are per calculation. All methods give identical $S,A,D,\ell$ and final grids. Compilation (\R{compileSeconds}\,s) and initial loading/warm-up (\R{warmupSeconds}\,s) are timed separately. Ranges are minimum to maximum. Speed ratios use the Python scan median. Times exclude generating input sites and reading or writing data.}
\label{tab:performance}
\end{table}

The Python queue is \R{queueSpeed} times faster than the scan, combining active-site selection, consecutive topplings and stored neighbour lists. Processing the full input sequence in C further reduces the overhead of calls between Python and C. The environment used Python 3.12.2, NumPy 1.26.4, SciPy 1.13.0, Matplotlib 3.8.4 and Apple Clang 17 (arm64 macOS, eight logical CPUs).

\section{Simulation lengths and initial conditions}\label{sec:design}
The preliminary runs used $L=64$ for open uniform, open biased and cylindrical input, each starting from empty, random stable and all-three grids with independent random-number sequences. They followed $60L^2$ grains in the uniform and cylindrical cases and $100L^2$ under bias. A separate batch comparison started from an equilibrated grid and checked 81,920 subsequent grains against the matched sequential run. Based on these runs, we used five seeds per condition, baseline sizes $L=8,16,32,64$, extension sizes $L=32,64$, and 200,000 measured grains per seed. Burn-in is $20L^2$ for baseline and cylinder, and $40L^2$ for bias. Batches start from the matched baseline's grid after burn-in.

The 50 primary runs contain \R{productionGrains} measured grains and \R{productionEvents} events, plus 2,080,000 burn-in grains. Each batch run has 50,000 events and uses the same initial grid and subsequent input sites as its single-grain control. Those two million batch grains form paired comparisons. Separate random-number sequences are used for the pilot, primary conditions apart from matched batches, toppling-order checks, the 100 independent $L=32$ ensemble chains and block resampling. The later confirmation experiment in Section~\ref{sec:bias} uses additional seeds and grains. The primary simulations took \R{productionWall}\,s, including \R{kernelSeconds}\,s in simulation calculations and \R{ioSeconds}\,s in data storage. These times include burn-in.

\section{Uncertainty and distribution fitting}\label{sec:statistics}
Primary 95\% Student-$t$ intervals use the variation among five independent run estimates, with four degrees of freedom. This allows events within each run to be correlated. Independent extension differences use Welch intervals. Batch comparisons use exact pairing. The time-series comparison uses 24 complete contiguous blocks of 8,192 grains, covering 196,608 grains. The remaining grains are retained in all-event summaries. Block lengths of 4,096 and 16,384 grains and 200 block-bootstrap resamples assess sensitivity. Autocorrelation is calculated separately for height, topplings, loss, dissipation indicators and height frequencies. The initial-positive-pair estimator uses a finite lag range and is bounded below by one, so it can return one for short-lag anticorrelation.

Fits use the exact counts of raw integer observations, with normalized likelihood
\begin{equation}
 p(x\mid\alpha,\lambda)=\frac{x^{-\alpha}e^{-\lambda x}}
 {\sum_{n=x_{\min}}^{x_{\max}}n^{-\alpha}e^{-\lambda n}}.
 \label{eq:fit}
\end{equation}
The candidates are power ($\lambda=0$), cutoff power and exponential ($\alpha=0$), fitted to identical observations and support. Fitted free parameters lie inside the bounds $-10\leq\alpha\leq10$ and $0\leq\lambda x_{\max}\leq10^5$. Finite support allows negative effective slopes. Each seed must retain at least 1,000 observations. The fixed lower limits are $(16,8,1,3)$ for $(S,A,D,\ell)$. Sensitivity checks halve and double them, with a minimum of one. For $S$, the upper fitting limit $L^2$ selects an intermediate range. The other upper limits, $L^2$, $3L^2+1$ and $2L-1$, are physical bounds. We record observed maxima separately from these fitting limits.

Parameter intervals are calculated across the five seed fits. Kolmogorov--Smirnov (KS) distances measure discrepancies between fitted and empirical CDFs. The cutoff model has one additional parameter. KS distances and marginal log scores provide descriptive comparisons on the observed sample, with block resampling used to assess the log-score comparison under temporal dependence. Logarithmically binned PMFs divide counts by bin width and sample count. Fitting uses the original integer observations, with each size and model analysed separately.
